\section{Introduction} \label{sec:1}

The detection of gravitational waves represents one of the most profound experimental achievements
in the history of physics. Predicted by Albert Einstein's general theory of relativity in 1916,
gravitational waves are ripples in the fabric of spacetime produced by the acceleration of massive
objects. For nearly a century, they remained beyond observational reach --- their amplitudes so
vanishingly small that detecting them demanded sensitivities previously thought impossible. That
barrier was broken on 14 September 2015, when the Laser Interferometer Gravitational-Wave
Observatory (LIGO) recorded GW150914, the first direct observation of a gravitational-wave signal,
originating from the inspiral and merger of two stellar-mass black holes approximately 1.3 billion
light-years away \cite{Abbott2016Discovery}. The event marked not merely a technological milestone but a
turning point in observational astronomy: it confirmed the existence of binary black hole systems,
validated general relativity in the strong-field dynamical regime, and opened an entirely new
observational window onto the universe, one inaccessible to electromagnetic astronomy.

\subsection*{Gravitational Waves from Linearized General Relativity}

Gravitational waves emerge naturally from Einstein's field equations when spacetime is treated as
a flat Minkowski background perturbed by a small metric deviation. Writing the full metric as
\begin{equation}
	g_{\mu\nu} = \eta_{\mu\nu} + h_{\mu\nu}, \qquad |h_{\mu\nu}| \ll 1,
\end{equation}
where $\eta_{\mu\nu} = \mathrm{diag}(-1,+1,+1,+1)$ is the flat Minkowski metric and $h_{\mu\nu}$
is the perturbation, and introducing the trace-reversed perturbation
\begin{equation}
	\bar{h}_{\mu\nu} = h_{\mu\nu} - \frac{1}{2}\,\eta_{\mu\nu}\, h, \qquad
	h \equiv \eta^{\mu\nu}h_{\mu\nu},
\end{equation}
the Lorenz gauge condition $\partial^{\mu}\bar{h}_{\mu\nu} = 0$ reduces the linearized Einstein
equations to the compact wave equation
\begin{equation}
	\Box\,\bar{h}_{\mu\nu} = -\frac{16\pi G}{c^{4}}\,T_{\mu\nu},
	\label{eq:wave}
\end{equation}
where $\Box = -c^{-2}\partial_t^2 + \nabla^2$ is the d'Alembertian operator and $T_{\mu\nu}$ is
the stress-energy tensor of the source. In vacuum ($T_{\mu\nu}=0$), Eq.~\eqref{eq:wave} becomes
the homogeneous wave equation, admitting plane-wave solutions propagating at the speed of light.
In the transverse-traceless (TT) gauge, only two independent polarization states survive --- the
\emph{plus} polarization $h_{+}$ and the \emph{cross} polarization $h_{\times}$ --- reflecting the
two physical degrees of freedom of a massless spin-2 field. A passing gravitational wave with
wavevector along the $z$-axis deforms a ring of test masses in the $xy$-plane as
\begin{equation}
	ds^2 = -c^2\,dt^2 + \bigl[1 + h_+(t)\bigr]\,dx^2
	+ \bigl[1 - h_+(t)\bigr]\,dy^2
	+ 2\,h_\times(t)\,dx\,dy + dz^2,
\end{equation}
making interferometric arm-length differences directly sensitive to the gravitational-wave strain
$h(t) = \Delta L(t)/L$.

\subsection*{The Quadrupole Formula and Binary Black Hole Systems}

Far from the source, the leading-order gravitational-wave emission is governed by the quadrupole
formula. The retarded solution to Eq.~\eqref{eq:wave} gives, in the TT gauge,
\begin{equation}
	h_{ij}^{\mathrm{TT}}(t,\mathbf{x}) =
	\frac{2G}{c^{4}\,r}\,\Lambda_{ij,kl}\,
	\ddot{Q}^{kl}\!\left(t_{\mathrm{ret}}\right), \qquad t_{\mathrm{ret}} = t - r/c,
	\label{eq:quad}
\end{equation}
where $r$ is the luminosity distance to the source, $\Lambda_{ij,kl}$ is the TT projection
operator, and $Q^{ij}$ is the traceless mass quadrupole moment\cite{Maggiore2007}
\begin{equation}
	Q^{ij} = \int_{\mathcal{V}} \rho\!\left(x^i x^j - \tfrac{1}{3}\,\delta^{ij} r^2\right)d^3x.
\end{equation}
Equation~\eqref{eq:quad} encodes a fundamental selection rule: gravitational-wave emission
requires a \emph{time-varying quadrupole moment}. Spherically symmetric collapse and uniformly
rotating rigid bodies produce no radiation; a binary system, by contrast, is an ideally
asymmetric and time-varying source.

For two point masses $m_1$ and $m_2$ in a circular Newtonian orbit of separation $d$ and
angular frequency $\Omega$, the total radiated power follows the Peters formula
\begin{equation}
	P = \frac{32}{5}\,\frac{G^4}{c^5}\,
	\frac{(m_1 m_2)^2(m_1+m_2)}{d^5}.
	\label{eq:peters}
\end{equation}
Energy loss through gravitational radiation causes the orbital separation to shrink and the
orbital frequency to increase. The merger timescale for a circular binary is
\begin{equation}
	\tau_{\mathrm{merge}} \approx
	\frac{12}{19}\,\frac{c^5\,d_0^4}{G^3\,m_1\,m_2\,(m_1+m_2)},
\end{equation}
where $d_0$ is the initial separation \cite{Maggiore2007}. For the GW150914 progenitor system this timescale was on
the order of the age of the universe, culminating in the observed merger event.

\subsection*{The Chirp Signal and the Chirp Mass}

The frequency evolution of the gravitational-wave signal during inspiral is governed by
\begin{equation}
	\dot{f}_{\mathrm{GW}} =
	\frac{96\pi^{8/3}}{5}\,
	\left(\frac{G\mathcal{M}}{c^3}\right)^{5/3}
	f_{\mathrm{GW}}^{\,11/3},
	\label{eq:fdot}
\end{equation}
where the \emph{chirp mass}
\begin{equation}
	\mathcal{M} = \frac{(m_1 m_2)^{3/5}}{(m_1+m_2)^{1/5}}
\end{equation}
is the single best-constrained quantity in a gravitational-wave observation.
Integrating Eq.~\eqref{eq:fdot} shows that $f_{\mathrm{GW}}$ increases monotonically with time,
and the waveform amplitude scales as $h \propto \mathcal{M}^{5/3} f_{\mathrm{GW}}^{2/3} / r$ \cite{Cutler_1994}.
The resulting signal --- a sinusoid with simultaneously increasing frequency and amplitude --- is
universally known as a \emph{chirp}. For GW150914, with $m_1 \approx 36\,M_\odot$ and
$m_2 \approx 29\,M_\odot$, the chirp mass is $\mathcal{M} \approx 28.3\,M_\odot$, and the
signal sweeps from $\sim 35$\,Hz to $\sim 250$\,Hz in roughly $0.2$ seconds \cite{Abbott2016Discovery}.

\subsection*{Three Phases of the Signal}

The complete gravitational-wave signal from a binary black hole coalescence divides into three
morphologically distinct phases:

\begin{enumerate}
	\item \textbf{Inspiral.} The early phase is well described by post-Newtonian (PN) perturbation
	theory, in which the binary slowly loses binding energy to gravitational radiation. The
	waveform is a chirp with monotonically increasing frequency and amplitude, whose phase
	evolution is captured order-by-order in the small parameter $v/c$. For GW150914, the
	inspiral sweeps through the LIGO band from $\sim 35$\,Hz to $\sim 150$\,Hz \cite{Cutler_1994, Maggiore2007}.
	
	\item \textbf{Merger.} As the separation approaches the light-ring radius, post-Newtonian
	expansions diverge and the system enters the fully nonlinear regime of general relativity.
	This phase, characterized by peak strain amplitude and highly asymmetric mass distribution,
	is accessible only through numerical relativity simulations, such as those underpinning the
	IMRPhenomD family of waveform models used in the present analysis\cite{Cutler_1994, Maggiore2007}.
	
	\item \textbf{Ringdown.} Following merger, the remnant Kerr black hole relaxes to stationarity
	by emitting gravitational radiation in the form of exponentially damped sinusoids known as
	quasi-normal modes (QNMs). The dominant QNM frequency and quality factor are determined
	entirely by the remnant mass $M_f$ and dimensionless spin $\chi_f$, providing a direct
	observational test of the Kerr no-hair theorem \cite{Cutler_1994, Maggiore2007}.
\end{enumerate}

\subsection*{Scope of This Work}

This paper presents an independent analysis of GW150914 using publicly available strain data from
the LIGO Open Science Center. The goal is to reproduce and interpret the key observational
signatures of the event through a sequence of complementary methods. We begin with data
acquisition and conditioning --- including power spectral density estimation, whitening, and
bandpass filtering --- to suppress instrumental noise and isolate the astrophysical signal
(Section~3). We then measure the inter-detector time delay between the Hanford (H1) and
Livingston (L1) observatories via cross-correlation, providing direct evidence of signal coherence
across the detector network (Section~4). A time--frequency analysis using the Q-transform is
employed to visualize the characteristic chirp morphology (Section~5). Finally, we apply matched
filtering with IMRPhenomD templates to quantify the signal-to-noise ratio, and perform a
$\chi^2$ signal consistency test to assess template match quality and suppress noise transients
(Section~6).

The combination of these analyses yields a reweighted SNR $\hat{\rho} \approx 13.3$, a reduced
$\chi^2_{\mathrm{red}} = 2.10$, and best-fit component masses of $35.5\,M_\odot$ and
$30.5\,M_\odot$, in close agreement with the published discovery paper \cite{Abbott2016Discovery}. Taken
together, they constitute a self-contained and reproducible confirmation of the first
gravitational-wave detection, and illustrate the power of open data in enabling independent
scientific verification.